
\subsubsection{2.1 Group-Robust Training}

The canonical approach to worst-group accuracy degradation is Group Distributionally Robust Optimization (Group DRO) \citep{Sagawa2019}, which minimizes the worst-case expected loss across predefined groups. Group DRO and its variants achieve 84--91\% WGA on Waterbirds and CelebA --- but require group annotations at training time. G2-SAM \citep{park2025g2sam} extends this to SAM-based optimization, applying group-wise perturbation to reduce per-group sharpness in supervised settings. G2-SAM demonstrates that sharpness reduction improves worst-group accuracy --- the geometric evidence that motivates our SSL extension. However, G2-SAM requires explicit group labels and operates in supervised cross-entropy settings; we remove both constraints.

Domain-Generalization SAM (DGSAM) \citep{Song2025} applies SAM variants to domain generalization rather than within-distribution subgroup robustness; our setting is orthogonal.

\subsubsection{2.2 Annotation-Free Debiasing}

Several methods achieve group robustness without group labels. Just Train Twice (JTT) \citep{liu2021just} trains an initial model, identifies misclassified samples as likely minority examples, and upweights them in a second training run. Loss-based Feature Resampling (LFR) \citep{ghaznavi2023lfr} uses high-loss samples from an initial linear probe as a proxy for minority group membership, enabling annotation-free resampling that matches oracle-labeled performance on Waterbirds and CelebA. EVaLS \citep{ghaznavi2024evals} extends LFR with environment inference for model selection. Environment Inference for Invariant Learning (EIIL) \citep{Creager2021} infers environment structure from gradient statistics.

All of these methods operate post-hoc on fixed representations: they take an SSL or ERM model as given and apply downstream reweighting. They cannot change the geometric structure of the representation itself. Our approach instead intervenes during SSL pretraining --- the geometric formation phase --- to produce representations with lower spurious anisotropy from the start.

\subsubsection{2.3 SSL and Spurious Correlations}

Izmailov et al. [2022] demonstrated that SSL models (SimCLR, DINO, Barlow Twins) can achieve competitive worst-group accuracy on Waterbirds and CelebA with appropriate feature retraining (DFR), suggesting SSL representations contain the necessary information for core-feature classification. Zhang \& Ré [2022] documented that CLIP models exhibit severe spurious correlation sensitivity despite strong average performance --- a pattern consistent with our DINO findings. Cross-Variant SSL \citep{yadav2026crossvariant} addresses shortcut reliance by diversifying the augmentation pipeline with generative models, achieving 92.5\% WGA on Waterbirds [UNVERIFIED --- verify before submission]. This approach modifies what the SSL model sees (augmentation diversity) rather than how it optimizes (loss landscape geometry). Our work is complementary: augmentation changes the data distribution; SAM changes the curvature structure.

A recent NeurIPS 2025 paper by Chen et al. [2025] (''Mitigating Spurious Features in Contrastive Learning with Spectral Regularization'') analyzes the covariance singular modes of SimCLR, DINO, and BYOL representations on spurious benchmarks, and proposes a spectral regularization term that enforces a uniform eigenspectrum during pretraining. This work achieves 43.8\% WGA for SimCLR on Waterbirds and provides geometric (covariance-level) characterization of SSL spurious encoding --- complementary to our sharpness anisotropy approach. Note: the first author name for \texttt{chen2025spectral} must be verified against NeurIPS 2025 proceedings before submission.

\subsubsection{2.4 Loss Landscape Geometry and Simplicity Bias}

SAM \citep{foret2021sam} finds flat minima by minimizing the maximum loss within a perturbation ball. Flat minima generalize better empirically and connect theoretically to implicit regularization. Gatmiry et al. [2024] prove that SAM promotes rank-1 (simplicity) bias in supervised cross-entropy settings --- SAM in supervised learning tends to learn simpler (lower-rank) features first. Since spurious features are typically simpler (lower-rank) than core features, this creates a theoretical tension: does SAM increase or decrease shortcut reliance in SSL?

The Spectral Curvature-Embedding Representation (SCER) framework \citep{park2025g2sam} theoretically links embedding geometry to worst-group error in SSL settings, providing motivation for directional curvature analysis. However, SCER does not provide empirical measurements of sharpness anisotropy, and does not address optimizer-level interventions.

Chen et al. [2025] provide an objective-level intervention: modifying the InfoNCE loss to penalize dominant spurious singular modes in the feature covariance matrix. Our work differs in two key ways: (1) we intervene at the optimizer level rather than the objective level --- SAM requires no modification of the contrastive loss; (2) our geometric diagnostic is directional sharpness anisotropy (loss landscape curvature) rather than spectral covariance analysis. These are orthogonal geometric characterizations: covariance singular modes describe what the features encode; sharpness anisotropy describes how the loss landscape is curved.

We bridge these lines: we bring the geometric sharpness analysis of the SAM literature to the SSL spurious correlation setting, providing the first empirical measurement protocol and, to our knowledge, the first application of SAM during SSL pretraining for shortcut reduction. The critical theoretical question --- whether Gatmiry's simplicity bias transfers to InfoNCE objectives --- is the central empirical question this work addresses.

\begin{table*}[htbp]
\centering
\resizebox{\dimexpr 2\columnwidth+\columnsep\relax}{!}{%
\begin{tabular}{lllll}
\toprule
Method & Annotation-Free & SSL Pretraining & Geometric Diagnosis & Training-Time \\
\midrule
Group DRO [Sagawa 2019] & No & No & No & Yes \\
JTT [Liu 2021] & Yes & No & No & Yes (2-stage) \\
LFR [Ghaznavi 2023] & Yes & No & No & Post-hoc \\
G2-SAM [Park 2025] & No & No & Yes (sharpness) & Yes \\
Cross-Variant SSL [Yadav 2026] & Yes & Yes & No & Yes (augment) \\
Chen et al. [2025] spectral-reg & Yes & Yes & Yes (covariance) & Yes (objective) \\
**Ours** & **Yes** & **Yes** & **Yes (anisotropy)** & **Yes (optimizer)** \\
\bottomrule
\end{tabular}
}
\end{table*}

